\section{Conclusion}
\label{sec:conclusion}

County preservation is a constitutionally required criterion in 34 state
redistricting frameworks.  Existing algorithmic approaches either ignore it
(producing unnecessary splits) or impose hard constraints (producing
infeasible instances).  We introduced county-sticky edge weighting as a
practical middle path.

The key findings are:

\begin{enumerate}
  \item \textbf{$\alpha_c = 3.0$ is Pareto-optimal}: The sweep of
    $\alpha_c \in \{1.0, 1.5, 2.0, 2.5, 3.0, 3.5, 5.0, 10.0\}$, including
    finer grid points around the elbow, robustly identifies $\alpha_c = 3.0$
    as the Pareto-frontier elbow.  Beyond $\alpha_c = 3.5$, the county-split
    reduction per unit of compactness cost drops by more than $4\times$.
  \item \textbf{Substantial split reduction}: At $\alpha_c = 3.0$,
    county splits fall by 34\% (487 $\to$ 323) and multi-county districts
    fall by 37\% (312 $\to$ 198) nationally.
  \item \textbf{Minimal compactness cost}: The mean Polsby--Popper score
    falls by only 3.0\% (0.361 $\to$ 0.350), leaving county-sticky maps
    well above the compactness of enacted congressional maps.
  \item \textbf{Population balance preserved}: Maximum population deviation
    increases by only 0.03 percentage points (well within the 0.5\%
    constitutional requirement).
  \item \textbf{State-level benefits scale with avoidable splits}: States
    with low district-to-county ratios (Iowa) see the largest reductions;
    large-population states (California) see minimal reduction because most
    splits are unavoidable.
\end{enumerate}

County-sticky weighting is implemented as \texttt{--weights-override county}
in the \texttt{redist} CLI and specified as the DIA default for states with
constitutional subdivision-preservation requirements.

\paragraph{Future work.}
Extensions include: (1) block-group resolution evaluation, (2) municipal
boundary preservation via a two-level hierarchy, (3) seed sensitivity
analysis under county-sticky weights, and (4) VRA-constrained county
preservation for majority-minority county clusters.
